% Standalone RevTeX appendix, or \input into a host manuscript after \appendix
% (omit everything from \documentclass through \begin{document} when using \input).
\documentclass[twocolumn,secnumarabic,amssymb,nobibnotes,aps,prd]{revtex4-2}

\usepackage{amsmath}
\usepackage{bm}
\usepackage{booktabs}

\begin{document}

\appendix

\section{Strong-to-Weak-Form Derivation for the Firedrake Solver}
\label{app:firedrake-weak-form}

This appendix documents the variational formulation implemented in
\texttt{solver-firedrake.py} for the coupled Saarelma--Connor pedestal
model~\cite{Saarelma2023}.
The three unknowns are flux-surface-averaged profiles on the radial
coordinate $x = r - r_{\mathrm{sep}}$:
\begin{equation}
  n_e(x),\quad \langle n_{\mathrm{FC}}\rangle(x),\quad \langle n_{\mathrm{CX}}\rangle(x),
\end{equation}
with $x=0$ at the separatrix and $x<0$ in the pedestal.
Coefficients $D_{\mathrm{ped}}(x)$, $S_i(x)$, $S_{\mathrm{CX}}(x)$,
$\langle|\nabla r|^{2}\rangle(x)$, and $|V_{\mathrm{CX}}|(x)$ are precomputed
on the parent grid and interpolated onto the Firedrake mesh;
$|V_{\mathrm{FC}}|$ and the form factors $f_{\mathrm{FC}}$, $f_{\mathrm{CX}}$
are treated as constants in the residual (currently $f_{\mathrm{FC}}=f_{\mathrm{CX}}=1$
in the Python code).

\subsection{Domain, function spaces, and boundary conditions}
\label{sec:domain}

The computational domain is the closed interval
\begin{equation}
  \Omega = [x_{\mathrm{inner}},\,0], \qquad x_{\mathrm{inner}} < 0.
\end{equation}
A one-dimensional \texttt{IntervalMesh} is used with boundary labels
\texttt{ds(1)} at $x=x_{\mathrm{inner}}$ (inner pedestal foot) and
\texttt{ds(2)} at $x=0$ (separatrix; inflow for neutrals).

The mixed finite-element space is
\begin{equation}
  W = V \times V \times V, \qquad
  u = \bigl(n_e,\,\langle n_{\mathrm{FC}}\rangle,\,\langle n_{\mathrm{CX}}\rangle\bigr) \in W,
\end{equation}
with $V = \mathrm{CG}_p(\Omega)$ (default $p=2$).
Test functions are denoted $v_e, v_F, v_C \in V$.

\paragraph{Boundary conditions.}
Four scalar conditions close the fourth-order coupled system:
\begin{subequations}
\label{eq:bcs}
\begin{align}
  n_e(0) &= n_{e,0}, \label{eq:bc-ne0}\\
  \langle n_{\mathrm{FC}}\rangle(0) &= \langle n_{\mathrm{FC}}\rangle_{0}, \label{eq:bc-nfc0}\\
  \langle n_{\mathrm{CX}}\rangle(0) &= \langle n_{\mathrm{CX}}\rangle_{0}, \label{eq:bc-ncx0}
\end{align}
\end{subequations}
with Dirichlet data at the separatrix, and at $x=x_{\mathrm{inner}}$ either
\begin{equation}
  n_e(x_{\mathrm{inner}}) = n_{e,\mathrm{inner}}
  \quad\text{(Dirichlet)}
  \label{eq:bc-ne-dirichlet}
\end{equation}
or
\begin{equation}
  \left.\frac{dn_e}{dx}\right|_{x=x_{\mathrm{inner}}} = n'_{e,\mathrm{inner}}
  \quad\text{(Neumann / flux)}.
  \label{eq:bc-ne-neumann}
\end{equation}
Neutrals have inflow at $x=0$ only; no boundary condition is imposed on
$\langle n_{\mathrm{FC}}\rangle$ or $\langle n_{\mathrm{CX}}\rangle$ at
$x=x_{\mathrm{inner}}$ (outflow).
Essential Dirichlet conditions are enforced with \texttt{DirichletBC}; the
Neumann condition for $n_e$ enters as a natural boundary term
(Sec.~\ref{sec:ne-weak}).

\subsection{Strong forms solved}
\label{sec:strong}

After flux-surface averaging and introduction of the form factors
(Eq.~(7) of Ref.~\cite{Saarelma2023}), the code solves the following system on
$\Omega$:
\begin{subequations}
\label{eq:strong-system}
\begin{align}
  \frac{d}{dx}\!\left[\langle|\nabla r|^{2}\rangle\,D_{\mathrm{ped}}\,\frac{dn_e}{dx}\right]
    &= -\,n_e\,S_i\,(\langle n_{\mathrm{FC}}\rangle + \langle n_{\mathrm{CX}}\rangle),
  \label{eq:strong-ne}\\[4pt]
  |V_{\mathrm{FC}}|\,\frac{d}{dx}\!\left[f_{\mathrm{FC}}\,\langle n_{\mathrm{FC}}\rangle\,\langle|\nabla r|^{2}\rangle\right]
    &= +\,n_e\,(S_i + S_{\mathrm{CX}})\,\langle n_{\mathrm{FC}}\rangle,
  \label{eq:strong-nfc}\\[4pt]
  |V_{\mathrm{CX}}|\,\frac{d}{dx}\!\left[f_{\mathrm{CX}}\,\langle n_{\mathrm{CX}}\rangle\,\langle|\nabla r|^{2}\rangle\right]
    &= n_e\,\left(\langle n_{\mathrm{CX}}\rangle\,S_i - \tfrac{1}{2}\,S_{\mathrm{CX}}\,\langle n_{\mathrm{FC}}\rangle\right).
  \label{eq:strong-ncx}
\end{align}
\end{subequations}
Equations~\eqref{eq:strong-nfc} and~\eqref{eq:strong-ncx} are written in
conservative form for the fluxes
$F_F \equiv |V_{\mathrm{FC}}|\,f_{\mathrm{FC}}\,\langle n_{\mathrm{FC}}\rangle\,\langle|\nabla r|^{2}\rangle$
and
$F_C \equiv |V_{\mathrm{CX}}|\,f_{\mathrm{CX}}\,\langle n_{\mathrm{CX}}\rangle\,\langle|\nabla r|^{2}\rangle$.
All quantities are in physical (SI) units; no non-dimensionalization
is applied in the Firedrake residual.

\subsection{Electron density: integration by parts}
\label{sec:ne-weak}

\paragraph{Weighted residual.}
Multiply Eq.~\eqref{eq:strong-ne} by a test function $v_e \in V$ and
integrate over $\Omega$:
\begin{equation}
  \int_{x_{\mathrm{inner}}}^{0}
    \frac{d}{dx}\!\left[\langle|\nabla r|^{2}\rangle\,D_{\mathrm{ped}}\,\frac{dn_e}{dx}\right] v_e \, dx
  = -\int_{x_{\mathrm{inner}}}^{0} n_e\,S_i\,(\langle n_{\mathrm{FC}}\rangle + \langle n_{\mathrm{CX}}\rangle)\, v_e \, dx.
  \label{eq:ne-weighted}
\end{equation}

\paragraph{Integration by parts.}
Define the diffusive flux
$q \equiv \langle|\nabla r|^{2}\rangle\,D_{\mathrm{ped}}\,n_e'$.
With primes denoting $d/dx$,
\begin{equation}
  \underbrace{\Bigl[\langle|\nabla r|^{2}\rangle\,D_{\mathrm{ped}}\,n_e'\, v_e\Bigr]_{x=x_{\mathrm{inner}}}^{x=0}}_{\text{boundary terms}}
  - \int_{x_{\mathrm{inner}}}^{0} \langle|\nabla r|^{2}\rangle\,D_{\mathrm{ped}}\,n_e'\, v_e' \, dx
  = -\int_{x_{\mathrm{inner}}}^{0} n_e\,S_i\,(\langle n_{\mathrm{FC}}\rangle + \langle n_{\mathrm{CX}}\rangle)\, v_e \, dx.
  \label{eq:ne-ibp}
\end{equation}

\paragraph{Boundary terms.}
At any boundary where $n_e$ is prescribed (Dirichlet), the test function
satisfies $v_e=0$.
At $x=0$, $n_e(0)=n_{e,0}$ is Dirichlet, so the boundary contribution at
$x=0$ vanishes.
At $x=x_{\mathrm{inner}}$ with Dirichlet data~\eqref{eq:bc-ne-dirichlet},
$v_e(x_{\mathrm{inner}})=0$ and the inner boundary term also vanishes.
At $x=x_{\mathrm{inner}}$ with Neumann data~\eqref{eq:bc-ne-neumann},
$v_e(x_{\mathrm{inner}})$ is not constrained and the prescribed gradient
enters through
\begin{equation}
  \Bigl[\langle|\nabla r|^{2}\rangle\,D_{\mathrm{ped}}\,n_e'\, v_e\Bigr]_{x=x_{\mathrm{inner}}}
  \longrightarrow
  \langle|\nabla r|^{2}\rangle(x_{\mathrm{inner}})\,D_{\mathrm{ped}}(x_{\mathrm{inner}})\,n'_{e,\mathrm{inner}}\, v_e(x_{\mathrm{inner}}).
  \label{eq:ne-neumann-natural}
\end{equation}

\paragraph{Volume weak form.}
Collecting terms, the semi-discrete residual for Eq.~\eqref{eq:strong-ne}
(without Picard lagging) is: find $n_e \in V$ such that, for all admissible
$v_e \in V$,
\begin{widetext}
\begin{equation}
  \int_{x_{\mathrm{inner}}}^{0} \langle|\nabla r|^{2}\rangle\,D_{\mathrm{ped}}\,n_e'\, v_e' \, dx
  - \int_{x_{\mathrm{inner}}}^{0} n_e\,S_i\,(\langle n_{\mathrm{FC}}\rangle + \langle n_{\mathrm{CX}}\rangle)\, v_e \, dx
  + \delta_{\mathrm{Neumann}}\,
    \langle|\nabla r|^{2}\rangle(x_{\mathrm{inner}})\,D_{\mathrm{ped}}(x_{\mathrm{inner}})\,n'_{e,\mathrm{inner}}\, v_e(x_{\mathrm{inner}})
  = 0,
  \label{eq:weak-ne}
\end{equation}
\end{widetext}
where $\delta_{\mathrm{Neumann}}=1$ for a Neumann inner BC and
$\delta_{\mathrm{Neumann}}=0$ for Dirichlet.
In Firedrake, the volume terms are
\texttt{(g\_fd * D\_fd * ne.dx(0) * v\_e.dx(0) - ne * Si\_fd * (nfc + nCX) * v\_e) * dx}
and the Neumann contribution is
\texttt{g\_fd * D\_fd * dne\_dx\_inner * v\_e * ds(1)}.

\subsection{Franck--Condon neutrals: strong-form Galerkin}
\label{sec:nfc-weak}

Equation~\eqref{eq:strong-nfc} is first order with characteristics
directed from the separatrix inward.
The implementation uses a strong-form Galerkin residual (no integration by
parts on the flux derivative), so that the discrete solution matches a
forward march from the inflow boundary $x=0$.
Multiplying by $v_F \in V$ and integrating gives
\begin{equation}
  \int_{x_{\mathrm{inner}}}^{0}
    \left\{
      |V_{\mathrm{FC}}|\,\frac{d}{dx}\!\left[f_{\mathrm{FC}}\,\langle n_{\mathrm{FC}}\rangle\,\langle|\nabla r|^{2}\rangle\right]
      - n_e\,(S_i + S_{\mathrm{CX}})\,\langle n_{\mathrm{FC}}\rangle
    \right\} v_F \, dx = 0,
  \label{eq:weak-nfc}
\end{equation}
for all $v_F \in V$.
Because $|V_{\mathrm{FC}}|$ and $f_{\mathrm{FC}}$ are spatially constant in the code,
\begin{equation}
  |V_{\mathrm{FC}}|\,\frac{d}{dx}\!\left[f_{\mathrm{FC}}\,\langle n_{\mathrm{FC}}\rangle\,\langle|\nabla r|^{2}\rangle\right]
  = \frac{d}{dx}\!\left[|V_{\mathrm{FC}}|\,f_{\mathrm{FC}}\,\langle|\nabla r|^{2}\rangle\,\langle n_{\mathrm{FC}}\rangle\right],
\end{equation}
which corresponds to \texttt{(VFC\_const * fFC\_const * g\_fd * nFC).dx(0) * v\_F * dx}.
Integration by parts would transfer the derivative onto $v_F$ and introduce
an inflow boundary term that does not preserve the upwind structure.

\subsection{Charge-exchange neutrals: strong-form Galerkin}
\label{sec:ncx-weak}

The same reasoning applies to Eq.~\eqref{eq:strong-ncx}:
\begin{equation}
  \int_{x_{\mathrm{inner}}}^{0}
    \left\{
      |V_{\mathrm{CX}}|\,\frac{d}{dx}\!\left[f_{\mathrm{CX}}\,\langle n_{\mathrm{CX}}\rangle\,\langle|\nabla r|^{2}\rangle\right]
      - n_e\,\left(\langle n_{\mathrm{CX}}\rangle\,S_i - \tfrac{1}{2}\,S_{\mathrm{CX}}\,\langle n_{\mathrm{FC}}\rangle\right)
    \right\} v_C \, dx = 0,
  \label{eq:weak-ncx}
\end{equation}
for all $v_C \in V$.
In Firedrake, $|V_{\mathrm{CX}}|$ is the spatially varying coefficient
\texttt{Vcx\_fd}; $f_{\mathrm{CX}}$ is constant.

\subsection{Coupled residual and Picard linearization}
\label{sec:coupled}

The total residual solved by Firedrake is
\begin{equation}
  F(u;\,v_e,v_F,v_C) = F_1 + F_2 + F_3 = 0,
  \label{eq:total-F}
\end{equation}
where $F_1$, $F_2$, $F_3$ are the forms derived above, plus the optional
Neumann boundary term in $F_1$.
Each nonlinear term is a product of two unknowns
(e.g.\ $n_e\,S_i\,\langle n_{\mathrm{FC}}\rangle$).
With \texttt{use\_picard=True}, one factor is lagged at iterate $u^{(k)}$
while solving for $u^{(k+1)}$:
$F_1$ lags $(\langle n_{\mathrm{FC}}\rangle+\langle n_{\mathrm{CX}}\rangle)$,
$F_2$ lags $n_e$, and $F_3$ lags $n_e$ and $\langle n_{\mathrm{FC}}\rangle$
in the CX source.
With \texttt{use\_picard=False}, all products use the current iterate and a
single SNES Newton solve is applied to the fully coupled residual.
Updates may be under-relaxed,
$u^{(k+1)} \leftarrow (1-\lambda)\,u^{(k)} + \lambda\,\tilde u^{(k+1)}$,
with $\lambda=\texttt{relax}$ (default $0.5$).

\subsection{Mapping to \texorpdfstring{\texttt{solver-firedrake.py}}{solver-firedrake.py}}
\label{sec:code-map}

\begin{table*}[t]
\caption{\label{tab:code-map}Correspondence between strong forms, variational residuals, and Firedrake symbols.}
\begin{ruledtabular}
\begin{tabular}{lll}
Strong form & Weak / Galerkin residual & Code symbol \\
\hline
Eq.~\eqref{eq:strong-ne} & Eq.~\eqref{eq:weak-ne} & \texttt{F1} (+ \texttt{F1\_neumann\_bc}) \\
Eq.~\eqref{eq:strong-nfc} & Eq.~\eqref{eq:weak-nfc} & \texttt{F2} \\
Eq.~\eqref{eq:strong-ncx} & Eq.~\eqref{eq:weak-ncx} & \texttt{F3} \\
\end{tabular}
\end{ruledtabular}
\end{table*}

Essential boundary conditions~\eqref{eq:bc-ne0}--\eqref{eq:bc-ncx0} and
optional Dirichlet data~\eqref{eq:bc-ne-dirichlet} are enforced via
\texttt{DirichletBC}.
The Neumann condition~\eqref{eq:ne-neumann-natural} is the only
$n_e$-related boundary term appearing explicitly in the residual.

\begin{thebibliography}{1}
\bibitem{Saarelma2023}
S.~Saarelma, R.~J.~Akers, B.~C.~Bray, W.~A.~Cooper, R.~J.~Goldston, P.~Helander, J.~H.~Harris, J.~W.~Hughes, A.~Kirk, A.~E.~Leonard \emph{et al.},
Nucl. Fusion \textbf{63}, 052002 (2023).
\end{thebibliography}

\end{document}
